\section{Sharp conditional memory with damping and adaptive commands}
\label{sec:r75-noisy-conditional}

The one-dependent family of Section~\ref{sec:r74-mixing-obstruction}
also admits a sharp positive theory under conditional damping.  The
extension below realizes the complete observed instrument, with one
physical dictionary shared by every command.  The classical sector
inequality and the scalar thresholds of Section~\ref{sec:r69-spectral}
are inherited ingredients, not new spectral-region results.  Earlier
noise-dependent process lower bounds, including an inverse-square-root
worst-phase dependence, are due to Fanizza, Lumbreras and Winter
\cite[Theorem 2 and Section 4.2]{Fanizza24}.  Here the lower and upper meet
for an explicit uniformly positive, exactly depolarizing instrument, and
the upper works simultaneously under adaptive rotation commands.

For a qubit set
\[
 U_\theta=\begin{pmatrix}
 \cos(\theta/2)&-\sin(\theta/2)\\
 \sin(\theta/2)&\cos(\theta/2)
 \end{pmatrix},\qquad
 E=\begin{pmatrix}3/4&0\\0&1/4\end{pmatrix},\qquad
 \tau=|0\rangle\langle0|,\qquad
 \mathcal D(X)=\tr(X)I/2.
\]
For $0\le\alpha\le1$ define
\begin{equation}\label{eq:r75-noisy-instrument}
 \begin{split}
 \Psi^{\alpha,\theta}_0(X)
   &=\tfrac18[\alpha U_\theta XU_\theta^*
                    +(1-\alpha)\mathcal D(X)],\\
 \Psi^{\alpha,\theta}_1(X)&=\tfrac14\tr(EX)\tau,\\
 \Psi^{\alpha,\theta}_2&=\mathcal D-\Psi^{\alpha,\theta}_0
                                      -\Psi^{\alpha,\theta}_1.
 \end{split}
\end{equation}
The initial state is $I/2$.  An allowed command specifies $\theta$; the
report is $s\in\{0,1,2\}$.  A command may be selected from the visible
history.  A finite known command alphabet is enough for every finite
lower witness below.  The same construction also gives Borel transition
tables for the entire angle circle; this last formulation counts labels,
not the encoding or exact sampling cost of a real-valued command.

This family does not satisfy the retained uniform-interior hypothesis
on every posterior, even when $\alpha<1$: report $1$ always prepares the
pure state $\tau$.  Repeated reports $1$ keep that pure state, so passing
to a fixed block length does not restore that hypothesis.  The finite
construction below instead uses the common geometry of the damped
rotation, the reset vertex and the remaining branch.

\begin{lemma}[Physical margins and predictive rank]
\label{lem:r75-noisy-physical}
For every $\alpha\in[0,1]$ and every angle, the maps
\eqref{eq:r75-noisy-instrument} form a quantum instrument, their sum is
exactly $\mathcal D$, and their report probabilities lie in
$\{1/8\}$, $[1/16,3/16]$ and $[11/16,13/16]$, respectively.
The prepared response is predictively stationary under each unobserved
command.  At a fixed angle its output process is stationary and
one-dependent.  When $\alpha>0$ and $\sin\theta\ne0$, its word Hankel rank
is three, and
\begin{equation}\label{eq:r75-noisy-cylinder}
 p(1\,0^n\,1)=\frac{8^{-n}}{64}
           \left(1+\frac{\alpha^n}{2}\cos(n\theta)\right)
 \qquad(n\ge0).
\end{equation}
At $\alpha=0$ the complete command response has rank two and is independent
of the angles.  For all $\alpha$, it has no one-label realization.
\end{lemma}
\begin{proof}
The first two Choi matrices are positive.  Their operator norms are at
most $1/4$ and $3/16$, while $J(\mathcal D)=I_4/2$.  Hence
$J(\Psi^{\alpha,\theta}_2)\succeq I_4/16$.  The asserted trace bounds
are unchanged by damping.  Since $\mathcal D(I/2)=I/2$, averaging any
unobserved command fixes the prepared predictive response, including
arbitrary subsequent commands and terminal POVMs.  At a fixed angle,
summing an intervening report applies $\mathcal D$ and therefore
factorizes the probabilities of the blocks on either side, exactly as in
Theorem~\ref{thm:r74-depolarizing-obstruction}.

Report $1$ has initial probability $1/8$ and prepares $\tau$.  Conditional
on $n$ zero reports its Bloch vector is
$\alpha^n(\cos(n\theta),\sin(n\theta))$, giving
\eqref{eq:r75-noisy-cylinder}.  All reached subnormalized states are real
symmetric, so the Hankel rank is at most three.  For $\alpha>0$ and
nonreal phase, the three modes $1/8$, $\alpha e^{i\theta}/8$ and
$\alpha e^{-i\theta}/8$ are distinct with nonzero coefficients.  The
$3\times3$ submatrix with entries $p(1\,0^{i+j}\,1)$ is therefore a
product of two nonsingular Vandermonde matrices and a nonsingular
diagonal matrix.  This proves rank three.
At $\alpha=0$ every reached state is diagonal, giving rank at most two;
the $2\times2$ cylinder minor has determinant $2^{-19}>0$.
Finally $p(1)=1/8$ and $p(11)=3/128\ne p(1)^2$ for every parameter,
excluding a one-label realization.
\end{proof}

Write $N_{\mathrm{mix}}(\alpha,\theta)$ for the least label count
reproducing every report word at one fixed angle, and define
$N_{\mathrm{pure}}$ with a counted pure start.  The same conclusions below
hold if every terminal qubit POVM is included: the uppers realize the
subnormalized quantum states themselves and the lowers use report words
alone.
Write $\mathcal K_M$ for the set of complex eigenvalues of $M\times M$
row-stochastic matrices, as in Section~\ref{sec:r69-spectral}.

\begin{lemma}[Relative cylinder obstruction and polygon realization]
\label{lem:r75-full-instrument-bounds}
For $\alpha>0$ and $\sin\theta\ne0$, every $M$-label mixed-start
realization satisfies
\begin{equation}\label{eq:r75-cylinder-spectrum}
 \alpha e^{i\theta}\in\mathcal K_M.
\end{equation}
Conversely, let $P_q$ be the regular $q$-gon with vertices
$e^{2\pi i j/q}$, $q\ge3$, one of which represents $\tau$.  If
$\alpha e^{i\theta}\in P_q$, its $q$ physical pure states form one
dictionary invariant under all three normalized branches, and give a
$q$-label mixed-start realization of the complete instrument.  A single
dictionary works for any collection of angles satisfying this condition.
\end{lemma}
\begin{proof}
For a candidate HMM put $v=\mu T_1$ and $d=T_1\one$, so that the cylinder
in \eqref{eq:r75-noisy-cylinder} is $vT_0^nd$.  Restrict $T_0$ to labels
reachable from $v$ and from which $d$ is eventually positive, obtaining
$S$.  For every pair $i,j$ in this productive set, fixed positive entry
and exit paths give constants $c_{ij}>0$ and integers $a_i,b_j\ge0$ with
\[
 c_{ij}(S^n)_{ij}\le p(1\,0^{a_i+n+b_j}\,1)\le C_{ij}8^{-n}.
\]
Thus $\rho(S)\le1/8$.  Cayley--Hamilton and the three distinct modes in
\eqref{eq:r75-noisy-cylinder} force
$\alpha e^{i\theta}/8$ to be an eigenvalue of $S$.
An irreducible diagonal block $B$ carrying it has spectral radius
$0<R\le1/8$.  The Perron normalization proved before
Proposition~\ref{prop:r70-relative-spectrum} gives a stochastic matrix
$Q$ with eigenvalue $\alpha e^{i\theta}/(8R)$ and eigenvector $u$.
For a stationary row $\pi$ of $Q$, $\pi u=0$ because this eigenvalue is
nonreal.  Consequently
\[
 Q'=8R Q+(1-8R)\one\pi
\]
is stochastic and has eigenvalue $\alpha e^{i\theta}$.  Its size is at
most $M$; padding by self-loop labels proves
\eqref{eq:r75-cylinder-spectrum}.  This normalizes only the irreducible
block carrying the mode, and does not assume a positive Perron vector for
a reducible productive core.

Identify the real Bloch plane with $\mathbb C$, taking $\tau$ to $1$.
Multiplication by $\alpha e^{i\theta}$ maps $P_q$ into itself whenever
that point is in $P_q$, by its invariance under multiplication by a
$q$th root.  This handles branch zero.  Branch one resets to the vertex
$1$.  If $r\in P_q$ and $h=\tr(\Psi_1(\sigma_r))\in[1/16,3/16]$, the
normalized third-branch Bloch vector is
\[
 \frac{-\alpha R_\theta r/8-h(1,0)}{7/8-h}.
\]
Its norm is at most $5/11$, strictly below the inradius
$\cos(\pi/q)\ge1/2$.  Thus it too lies in $P_q$.
Choose barycentric coefficients for each of these normalized images and
multiply them by the branch probability.  The resulting nonnegative
matrices $T_s$ have stochastic sum and satisfy, for every dictionary state,
\[
 \Psi_s(\sigma_i)=\sum_jT_s(i,j)\sigma_j.
\]
Uniform initial mixture represents $I/2$.  Induction proves this identity
for every subnormalized word state; terminal effects and common adaptive
policies therefore have the same laws.  Triangulation of the polygon
provides a piecewise continuous, hence Borel, choice when the angle varies.
\end{proof}

\begin{theorem}[Sharp noisy-instrument memory]
\label{thm:r75-noisy-memory}
For $M\ge3$, $0<\theta\le2\pi/M$ and $0<\alpha\le1$,
\begin{equation}\label{eq:r75-first-sector}
 N_{\mathrm{mix}}(\alpha,\theta)\le M
 \quad\Longleftrightarrow\quad
 \alpha\le\frac1{\cos\theta+\sin\theta\tan(\pi/M)}.
\end{equation}
At equality the mixed minimum is exactly $M$, and the pure minimum is
$M+1$.  In every $M$-label mixed realization at this boundary, its
zero-report matrix is diagonally similar to
$\frac18[(1-t)I+tC_M]$ for some $0<t\le1$, up to cyclic relabelling and
orientation.  No analogous claim is made for its other report matrices.

For $0<\alpha<1$ put
\begin{equation}\label{eq:r75-universal-count}
 K(\alpha)=\max\left\{3,
               \left\lceil\frac\pi{\arccos\alpha}\right\rceil\right\}.
\end{equation}
One dictionary of $K(\alpha)$ labels realizes simultaneously every
rotation command in \eqref{eq:r75-noisy-instrument}, with a mixed start.
This count is minimal for the full angle command family, and the
corresponding pure minimum is $K(\alpha)+1$.  In particular
\[
 \max_{\sin\theta\ne0}N_{\mathrm{mix}}(\alpha,\theta)=K(\alpha),
 \qquad
 \max_{\sin\theta\ne0}N_{\mathrm{pure}}(\alpha,\theta)=K(\alpha)+1.
\]
The same common upper holds for each finite known angle alphabet, with
matching lower whenever that alphabet contains a phase attaining the
displayed maximum.

At $\alpha=0$, the complete command response has exactly two labels with
a mixed start and three with a pure start.  At $\alpha=1$, the full angle
family has no finite realization; for one nonreal angle its finite-order
and infinite-order alternatives are those of
Theorem~\ref{thm:r74-depolarizing-obstruction}.  Whenever a relevant
minimum is finite, pure initialization adds exactly one label.
\end{theorem}
\begin{proof}
The lower in \eqref{eq:r75-first-sector} follows from
\eqref{eq:r75-cylinder-spectrum} and Lemma~\ref{lem:r69-sector}.
For the upper, the radial boundary of the edge between $1$ and
$e^{2\pi i/M}$ is the right side of
\eqref{eq:r75-first-sector}; use the polygon realization.  At that
boundary strict monotonicity of $\cot(\pi/M)$ excludes every smaller
size at least three; the rank-three lower excludes sizes one and two.
For any realizing size $M$, the irreducible block in the preceding
proof must use all $M$ labels.  Applying the sector inequality before
the final radial contraction gives
\[
 \frac{\alpha}{8R}
  [\cos\theta+\sin\theta\tan(\pi/M)]\le1,
\]
so $R\ge1/8$, and hence $R=1/8$.  Equality in the inherited sector lemma
therefore fixes its Perron normalization to the lazy directed cycle.
The stationary initialization theorem supplies the pure count without
requiring the original diagonal similarity to be scalar.

The disk of radius $\alpha$ lies in the regular $K(\alpha)$-gon, including
the equality case $\alpha=\cos(\pi/K(\alpha))$.  The preceding lemma
therefore constructs one physical dictionary for every angle.  If
$K=3$, any nonreal phase has rank three and requires three labels.  If
$K>3$, take $\theta=\pi/(K-1)$.  The definition of the ceiling gives
$\alpha>\cos(\pi/(K-1))$, so the first-sector criterion excludes
$K-1$ labels for this single phase.  It also excludes every smaller size
for the full command family.  The pure upper appends one initializer,
and the single critical phase, together with
Theorem~\ref{thm:r75-stationary-initializer}, gives the matching lower.

At $\alpha=0$, use the two diagonal density matrices with Bloch
coordinates $1$ and $-3/11$.  Their initial probabilities are $3/14$ and
$11/14$.  In that order the report matrices are
\begin{equation}\label{eq:r75-zero-damping-tables}
 T_0=\begin{pmatrix}3/112&11/112\\3/112&11/112\end{pmatrix},\qquad
 T_1=\begin{pmatrix}3/16&0\\19/176&0\end{pmatrix},\qquad
 T_2=\begin{pmatrix}0&11/16\\7/88&121/176\end{pmatrix}.
\end{equation}
Their sum has both rows $(3/14,11/14)$, and direct substitution verifies
the subnormalized state identities.  The rank-two lower and the
stationary initialization law give the asserted exact counts.
Finally the undamped full family includes a non-root phase, so its
infinite lower follows from the retained theorem.  For every finite
angle alphabet the prepared predictive state is fixed by every
unobserved command, and $p(11)\ne p(1)^2$; thus the exact initialization
law applies at all parameters, not just at spectral boundaries.
\end{proof}

For a fixed phase the lower in Lemma~\ref{lem:r75-full-instrument-bounds}
and its regular-polygon upper also transfer the already proved arithmetic
envelope of Theorem~\ref{thm:r71-envelope}, with the same constants, to
the complete observed instrument.  Consequently its two logarithmic
exponents and its badly-approximable criterion transfer as well.  This is
a change from a scalar response to an entire instrument, not a new
number-theoretic assertion.  No equality between the full-instrument
minimum and the scalar spectral minimum is asserted outside the proved
sectors and boundary cases.
